\documentclass[11pt]{article}
\usepackage[a4paper,margin=1in]{geometry}
\usepackage{amsmath,amssymb,amsthm,mathtools}
\usepackage{enumitem}

\newtheorem{theorem}{Theorem}[section]
\newtheorem{proposition}[theorem]{Proposition}
\newtheorem{lemma}[theorem]{Lemma}
\newtheorem{remark}[theorem]{Remark}

\newcommand{\R}{\mathbb R}
\newcommand{\Om}{\Omega}
\newcommand{\B}{B_1}
\newcommand{\Sph}{\mathbb S^{n-1}}
\newcommand{\eps}{\varepsilon}
\newcommand{\dT}{\delta_T}
\newcommand{\HH}{H^{1/2}}
\newcommand{\Lamb}{\Lambda}
\newcommand{\Th}{\theta}

\title{Why the Torsion Deficit Sees \(H^{1/2}\), Not \(H^1\),\\
and What the Capillary Obstruction Actually Rules Out}
\author{}
\date{May 18, 2026}

\begin{document}
\maketitle

\begin{abstract}
This note answers the question whether the Cauchy--Serrin defect \(D_H\),
together with the equation \(-\Delta u=1\), should give more than
\(H^{1/2}\)-smallness of a boundary perturbation.  The conclusion is:
boundary \(D_H\), if available in a controlled graph regime, is indeed an
\(H^1\)-scale quantity.  But the torsion deficit controls only an integrated
level-set \(D_H\), not the boundary \(D_H\).  High-frequency capillary ripples
can have small amplitude and prescribed \(H^{1/2}\) size while their torsion
deficit is much smaller than the square of that \(H^{1/2}\) size.  Thus a
bare \(H^{1/2}\)-coercive perturbative theorem is false.  This does not refute
the main sharp stability theorem, whose target is
\(\dT(\Om)\gtrsim\mathcal A(\Om)^2\), since Fraenkel asymmetry is an
\(L^1\)-scale quantity rather than an \(H^{1/2}\)-scale quantity.
\end{abstract}

\section{Three different quantities}

Let \(u_\Om\) solve
\[
   -\Delta u_\Om=1,\qquad u_\Om\in H^1_0(\Om),
\]
and write \(T(\Om)=\int_\Om u_\Om\).  At fixed volume the ball maximizes
torsion, and the Saint--Venant deficit is comparable to
\[
   \dT(\Om)=T(\B)-T(\Om)
\]
up to the harmless sign and factor conventions used elsewhere in the project.

For a smooth near-ball
\[
   \partial\Om_\eps=\{(1+\eps\varphi(\theta))\theta:\theta\in\Sph\},
\]
with volume and barycenter normalized, the first shape derivative satisfies
\[
   \Delta u_1=0\quad\hbox{in }\B,\qquad
   u_1|_{\partial\B}=\frac{\varphi}{n}.
\]
If \(\varphi=\sum_{k\ge0}\varphi_k\) is the spherical-harmonic decomposition,
then
\[
   u_1(r,\theta)=\frac1n\sum_{k\ge0}r^k\varphi_k(\theta),
   \qquad
   \int_{\B}|\nabla(nu_1)|^2
      =\sum_{k\ge1} k\,\|\varphi_k\|_{L^2(\Sph)}^2.
\]
This is exactly the \(H^{1/2}\) scale.

On the other hand, the boundary normal derivative has first variation
\[
   |\nabla u_{\Om_\eps}|_{\partial\Om_\eps}
   =\frac1n-\frac{\eps}{n}\sum_{k\ge0}(k-1)\varphi_k+O(\eps^2)
\]
in the classical small \(C^2\) graph regime.  Therefore a boundary Serrin or
Cauchy oscillation defect has the scale
\[
   D_H^{\partial}(\Om_\eps)
      \simeq \eps^2\sum_{k\ge2}(k-1)^2
             \|\varphi_k\|_{L^2(\Sph)}^2,
\]
which is an \(H^1\)-scale quantity.  The perimeter deficit has the same high
frequency scale,
\[
   P(\Om_\eps)-P(\B)\simeq
   \eps^2\sum_{k\ge2}k^2\|\varphi_k\|_{L^2(\Sph)}^2.
\]

Thus there is no contradiction:
\[
   \hbox{torsion deficit}\sim H^{1/2},\qquad
   \hbox{boundary }D_H\hbox{ or perimeter}\sim H^1.
\]
The issue is that the sharp stability problem gives the former, not the latter.

\section{Why the global \(D_H\) control does not give boundary \(H^1\)}

The level-set identity used in Plan 2 controls an integral of the form
\[
   \int \bigl(D_H(t)+D_I(t)\bigr)\,w(t)\,dt
      \le C\,\dT(\Om),
\]
or, in volume-radius variables, an analogous \(d\mu=(-t'(\rho))\,d\rho\)
integral over level surfaces.  This is an averaged statement over interior
levels.  It is not a pointwise bound for the boundary level \(t=0\).

In a coherent smooth foliation, \(D_H\) supplies the radial kinetic part of an
annular Sobolev norm and \(D_I\) supplies the angular part.  A trace from a
bulk \(H^1\)-type annular norm to one endpoint gives exactly an
\(H^{1/2}\)-boundary norm.  This is the same half-derivative that appears in
the Steklov spectrum of the ball.

To obtain \(H^1\)-smallness of the boundary graph one would need substantially
more information, for example a pointwise boundary \(D_H^\partial\) bound, a
uniform nondegenerate graph collar, or a theorem excluding concentration of the
level-set \(D_H\) in a very thin boundary layer.  The equation
\(-\Delta u=1\) does smooth the solution inside the domain, but without
boundary geometry it does not prevent the boundary from oscillating in a
region of tiny harmonic or torsional relevance.

\section{A capillary counterexample to bare \(H^{1/2}\) coercivity}

The too-strong low-regularity perturbative theorem would say that, for radial
graphs with small amplitude and small \(H^{1/2}\) norm,
\[
   \dT(\Om_\varphi)\ge c\,\|\varphi\|_{H^{1/2}(\Sph)}^2
\]
with no slope, perimeter, curvature, or higher-norm hypothesis.  This is
false already in dimension two.

\begin{proposition}[High-frequency capillary ripples]\label{prop:capillary}
Let \(n=2\).  For \(k\ge2\) and \(0<a\ll1\), set
\[
   R_{a,k}(\theta)=c_a(1+a\cos(k\theta)),\qquad
   \Om_{a,k}=\{re^{i\theta}:0<r<R_{a,k}(\theta)\},
\]
where \(c_a\) is chosen so that \(|\Om_{a,k}|=|B_1|\).  Then the barycenter of
\(\Om_{a,k}\) is zero,
\[
   \|R_{a,k}-1\|_{L^\infty}\le C a,\qquad
   \|R_{a,k}-1\|_{H^{1/2}(S^1)}^2
      = c\,k a^2+O(a^2+a^4),
\]
but
\[
   0\le T(B_1)-T(\Om_{a,k})\le C a .
\]
Consequently there are smooth volume-normalized, barycentered radial graphs
with \(\|\varphi_j\|_{L^\infty}\to0\),
\(\|\varphi_j\|_{H^{1/2}}^2\to0\), and
\[
   T(B_1)-T(\Om_j)
      =o\bigl(\|\varphi_j\|_{H^{1/2}(S^1)}^2\bigr).
\]
\end{proposition}

\begin{proof}
Since
\[
   \frac1{2\pi}\int_0^{2\pi}(1+a\cos(k\theta))^2\,d\theta
      =1+\frac{a^2}{2},
\]
the volume normalization is achieved by
\[
   c_a=(1+a^2/2)^{-1/2}=1+O(a^2).
\]
Thus \(\|R_{a,k}-1\|_{L^\infty}\le Ca\), and the constant Fourier mode is
only \(O(a^2)\).  The barycenter is proportional to
\[
   \int_0^{2\pi}R_{a,k}(\theta)^3 e^{i\theta}\,d\theta.
\]
For \(k\ge2\), the function \(R_{a,k}^3\) contains only Fourier modes that are
multiples of \(k\), hence this first Fourier coefficient vanishes.

The \(H^{1/2}(S^1)\) estimate follows from the Fourier characterization of
\(H^{1/2}\): the modes \(\pm k\) have amplitude comparable to \(a\), while
the constant correction contributes only \(O(a^4)\).  Hence
\[
   \|R_{a,k}-1\|_{H^{1/2}}^2=c\,k a^2+O(a^2+a^4).
\]

Finally,
\[
   B_{1-Ca}\subset \Om_{a,k}\subset B_{1+Ca}.
\]
Torsion is monotone under inclusion, and the ball maximizes torsion at fixed
volume.  Therefore
\[
   0\le T(B_1)-T(\Om_{a,k})
      \le T(B_1)-T(B_{1-Ca}).
\]
In two dimensions \(T(B_R)=\pi R^4/8\), so the last quantity is at most
\(Ca\).

Choose \(a_j=\eps_j^3\), \(k_j\simeq \eps_j^{-4}\), with \(\eps_j\to0\).
Then
\[
   \|R_{a_j,k_j}-1\|_{H^{1/2}}^2\simeq \eps_j^2,
   \qquad
   T(B_1)-T(\Om_{a_j,k_j})\le C\eps_j^3,
\]
which proves the claim.
\end{proof}

\begin{remark}
This does not contradict the classical second variation formula.  That formula
is local in a stronger topology, such as small \(C^{2,\gamma}\) or at least a
geometry that prevents the slope from blowing up.  In the example above
\(\|R_{a,k}'\|_{L^\infty}\sim ak\), which tends to infinity for the chosen
sequence.
\end{remark}

\section{What the equation \(-\Delta u=1\) can and cannot do}

The transformation
\[
   h(x)=u_\Om(x)+\frac{|x|^2}{2n}
\]
is useful: \(h\) is harmonic in \(\Om\), and the first variation at the ball is
the harmonic extension of the boundary displacement.  This is the reason the
natural spectral norm is \(H^{1/2}\).

However, harmonicity is an interior regularity statement.  It does not make
surface measure comparable to harmonic measure, and it does not force the
solution to see every oscillation of a rough boundary.  In the capillary
example, the oscillations live in a thin, high-frequency boundary layer.  The
torsion function can be very close to the ball torsion away from that layer,
and the cost of the layer is controlled by its amplitude, not by the
\(H^{1/2}\) size of the radial parametrization.

This is also the warning for two-dimensional conformal methods.  Dirichlet
energy is conformally natural, but the radial \(H^{1/2}\) norm in the
geometric angle \(\theta\) is not automatically comparable to the conformal
trace norm unless one has quantitative control of the boundary map
(for instance chord-arc, quasidisk, or quasisymmetry constants).  Without such
control, conformal coordinates compress the capillary fingers; the bad
oscillation may have very small harmonic measure.

\section{Compatibility with the main asymmetry theorem}

The previous proposition must not be read as a counterexample to the main
theorem.  In the same example,
\[
   |\Om_{a,k}\Delta B_1|
      =\frac12\int_0^{2\pi}|R_{a,k}(\theta)^2-1|\,d\theta
      = c\,a+O(a^2),
\]
with \(c>0\) independent of \(k\).  Hence
\[
   \mathcal A(\Om_{a,k})^2\simeq a^2.
\]
The upper bound \(T(B_1)-T(\Om_{a,k})\le Ca\) is completely compatible with
the sharp Saint--Venant conclusion
\[
   T(B_1)-T(\Om_{a,k})\ge c\,\mathcal A(\Om_{a,k})^2\simeq c\,a^2.
\]
What fails is only the stronger estimate with
\(\|R_{a,k}-1\|_{H^{1/2}}^2\simeq ka^2\) on the right.

\section{What extra ``tame capillary'' assumptions could work}

The obstruction above rules out bare \(H^{1/2}\)-norm coercivity, but not the
main theorem and not every perturbative route.  The following strengthened
statements remain plausible or already available.

\begin{enumerate}[label=\textbf{(\Alph*)}]
\item \textbf{Classical near-spherical regime.}
If \(\|\varphi\|_{C^{2,\gamma}}\) is sufficiently small, BDV's
near-spherical theorem gives
\[
   \dT(\Om_\varphi)\ge c(n)\|\varphi\|_{H^{1/2}}^2.
\]
Here the coercive norm is \(H^{1/2}\), but the smallness topology is much
stronger.

\item \textbf{Small \(H^{1/2}\) plus a fixed high-norm bound.}
If
\[
   \|\varphi\|_{H^{1/2}}\le\eps,\qquad
   \|\varphi\|_{C^{m,\gamma}}\le M(n,R)
\]
with \(m+\gamma>2+\gamma_0\), then interpolation gives
\[
   \|\varphi\|_{C^{2,\gamma_0}}
      \le C\,\|\varphi\|_{H^{1/2}}^\theta
              M(n,R)^{1-\theta}.
\]
For \(\eps\) small this enters BDV.  This is precisely the mechanism recorded
in the graph-bootstrap notes.

\item \textbf{Boundary \(D_H^\partial\) plus graph nondegeneracy.}
If the actual boundary Serrin defect is controlled and the boundary is already
a nondegenerate graph so that the expansion of \(|\nabla u|\) is valid, then
\[
   D_H^\partial\sim
   \|(\Lambda-I)\varphi\|_{L^2}^2
\]
does give \(H^1\)-scale control of the non-neutral modes.  The missing input in
the global problem is the boundary \(D_H^\partial\) bound itself.

\item \textbf{Uniform chord-arc, Lipschitz, NTA, or harmonic-measure control.}
In dimension two, a conformal or transport method could plausibly work under a
quantitative assumption preventing capillary screening, for example a uniform
quasidisk/chord-arc constant together with small amplitude.  Such a condition
would have to imply comparability between surface measure and harmonic measure
and a one-sided Dirichlet-to-Neumann form bound.  The constants would depend on
the tameness parameter.

\item \textbf{Interior levels.}
The equation \(-\Delta u=1\) smooths level sets away from the boundary.
Integrated \(D_H\) can provide radial kinetic control through the foliation,
and a trace theorem then naturally gives \(H^{1/2}\) control at a selected
level.  Transferring that information back to \(\partial\Om\) requires a
separate boundary-layer theorem excluding precisely the capillary behavior in
Proposition~\ref{prop:capillary}.
\end{enumerate}

\section{Conclusion}

The answer to the original question is therefore twofold.

First, \(D_H\) really is stronger than \(H^{1/2}\) when it is a boundary
Serrin defect in a controlled graph regime: spectrally it sees
\((\Lambda-I)\varphi\), hence \(H^1\).  But the sharp torsion deficit does not
control that boundary quantity.  It controls an integrated level-set defect,
and tracing an integrated \(H^1\)-type quantity to a boundary naturally loses a
half derivative.

Second, the equation \(-\Delta u=1\) is not enough by itself to recover the
missing half derivative at the rough boundary.  Capillary oscillations can be
screened from the torsion function.  Thus the perturbative theorem with
\(\|\varphi\|_{H^{1/2}}^2\) on the right is false under bare \(H^{1/2}\)
smallness.  The main theorem for this class has a different right hand side,
\(\mathcal A(\Om)^2\), and survives this example.  A Plan~3 proof that avoids
Plan~1 would have to prove an \(L^1/L^2\)-scale capillary estimate, such as a
lower bound by the radial second moment \(\int_{\Sph}(R-1)^2\), not a full
\(H^{1/2}\)-coercive estimate.

\end{document}
